\section{Related Works}
\label{sec:related}


\noindent\textbf{Text-to-Video Generation}.
Distinct from conventional unconditional and class-conditional video generation \cite{Wang2020ImaGINatorCS, Wang2019G3ANDA, Wang2021InMoDeGANIM, Chen2020LongTermVP, Clark2019AdversarialVG, Brooks2022GeneratingLV, Yu2022GeneratingVW, Skorokhodov2021StyleGANVAC, Tian2021AGI, Bhagat2020DisentanglingMF, Xie2019MotionBasedGM, Yan2021VideoGPTVG, Ge2022LongVG}, 
T2V generation focuses on automatically converting textual descriptions into videos. 
This is a challenging task,
as it involves understanding text semantics and translating it into video content.
This often requires 
powerful cross-modal algorithms \cite{radford2021learning}, 
large computing resources \cite{Hestness2017DeepLS, Chollet2019OnTM}, 
and extensive video data \cite{schuhmann2022laion5b, bain2021frozen, Wang2023LAVIEHV, Wang2022InternVideoGV, Xue2021AdvancingHV, Schuhmann2021LAION400MOD}.
Based on the success of diffusion models in the Text-to-Image (T2I) generation \cite{Nichol2021GLIDETP, rombach2022highResolution, Ramesh2022HierarchicalTI, saharia2022imagen, Balaji2022eDiffITD, Kusupati2022MatryoshkaRL, Dai2023EmuEI, BetkerImprovingIG},
a series of such works have been recently transferred to T2V generation \cite{singer2023makeavideo, Ho2022ImagenVH, blattmann2023align, Wang2023LAVIEHV, Ge2023PreserveYO, Zhang2023Show1MP, Wang2023VideoFactorySA, Wang2023ModelScopeTT}. 
%%%%%%%%%%%%%%%%%%%%
However,
whether trained from scratch \cite{Ho2022ImagenVH, singer2023makeavideo} or finetuned from T2I model \cite{blattmann2023align, Ge2023PreserveYO, Wang2023LAVIEHV, Wang2023ModelScopeTT, Zhang2023Show1MP, Wang2023VideoFactorySA}, 
most of these approaches mainly work on generating short videos with few seconds from simple descriptions.
In contrast,
our Vlogger can generate a minute-long vlog with complex stories.



\noindent\textbf{Long Video Generation}.
The generation of long videos predominantly relies on parallel \cite{Yin2023NUWAXLDO} or autoregressive \cite{villegas2023phenaki} structures.
However,
these early works still face challenges for vlog generation.
On the one hand,
the parallel manner can relax spatial-temporal content incoherence problems by coarse-to-fine generation.
However, this approach necessitates extensive and laborious training on large, well-annotated long video datasets \cite{Yin2023NUWAXLDO}.
On the other hand,
the autoregressive manner can relax heavy data requirements for training long videos,
by applying short video generation models iteratively with sliding windows.
However,
this solution often suffers from noticeable shot change and long-term incoherence \cite{villegas2023phenaki, chen2023seine, he2023latent, Qiu2023FreeNoiseTL, Wang2023GenLVideoMT}
which becomes problematic when generating vlogs encompassing complex narratives and multiple scenes. 
It is worth mentioning that the community has gradually realized that delegating higher-order reasoning tasks to LLM is of great help to visual tasks \cite{Wu2023VisualCT, gupta2023visual, Li2023VideoChatCV, Zhang2023ControllableTG, Liu2023VisualIT, Zhu2023MiniGPT4EV, Huang2023FreeBloomZT, BetkerImprovingIG}.
Our Vlogger introduces LLM to the field of long video generation,
and effectively addresses the problems of training burden and content incoherence in the previous methods,
by distinct cooperation of top-down planning and bottom-up shooting.
